% Header: General study / work / project template
%%%%%%%%%%%%%%%%%%%%%%%%%%%%%%%%%%%%%%%%%%%%%%%%%%%%%%%%%%%%%%%%%%%

\documentclass[11pt,a4paper]{scrartcl}

% Layout
\usepackage[margin=2.5cm]{geometry}
\usepackage[onehalfspacing]{setspace}

% General packages
\usepackage{graphicx}
\usepackage{enumitem}
\usepackage{tcolorbox}
\tcbuselibrary{breakable}
\usepackage[breaklinks=true,colorlinks=true,linkcolor=blue,urlcolor=blue,citecolor=blue]{hyperref}
\usepackage{amsmath,amsthm,amssymb,mathtools}
\usepackage{icomma}
\usepackage[T1]{fontenc}
\usepackage[utf8]{inputenc}

% Language: choose one
% \usepackage[ngerman]{babel}
\usepackage[english]{babel}

% Units / tables / floats / references
\usepackage[locale = UK,space-before-unit=true,per-mode = symbol]{siunitx}
\usepackage{booktabs,multirow}
\usepackage{placeins}
\usepackage[natbib,abbreviate=true,doi=false,style=numeric-comp,giveninits=true,sorting=none]{biblatex}
\usepackage{csquotes}

% Optional bibliography
% \addbibresource{references.bib}

% Graphics paths
\graphicspath{{Bilder/} {images/} {figures/} {chapters/}}

% Optional math shortcuts
\newcommand{\R}{\mathbb{R}}
\newcommand{\N}{\mathbb{N}}
\newcommand{\Q}{\mathbb{Q}}
\newcommand{\set}[1]{\left\{ #1 \right\}}
\newcommand{\abs}[1]{\left| #1 \right|}
\newcommand{\norm}[1]{\left\lVert #1 \right\rVert}
\newcommand{\ol}[1]{\overline{#1}}

% Optional units
\DeclareSIUnit{\dBm}{dBm}

% Box styles
\newtcolorbox{calculationbox}[1][]{
    title=Calculation,
    breakable,
    #1
}

\newtcolorbox{solutionbox}[1][]{
    title=Solution,
    breakable,
    #1
}

\newtcolorbox{answerbox}[1][]{
    title=Answer,
    breakable,
    #1
}

\newtcolorbox{notebox}[1][]{
    title=Notes,
    breakable,
    #1
}

\newtcolorbox{sidecalcbox}[1][]{
    title=Side Calculation,
    breakable,
    #1
}

%%%%%%%%%%%%%%%%%%%%%%%%%%%%%%%%%%%%%%%%%%%%%%%%%%%%%%%%%%%%%%%%%%%
\begin{document}

    \title{Physik 4 - Blatt 05}
    \author{Tobias Laser}
    \date{\today}
    \maketitle
    \vfill
    \thispagestyle{empty}

    \pagenumbering{arabic}

    \paragraph{Aufgabe 1: Q-Wert und CNO-Zyklus}

    Für die Energiebilanz einer Kernreaktion $A + a \rightarrow B + b$ wird der Q-Wert genutz. Dieser entspricht (bezugssystemunabhängig) der insgesamt freiwerdenden kinetischen Energie.

    Die astronomisch relevanteste Kernreaktion ist hierbei das Wasserstoffbrennen, welches in zwei Hauptreaktionen unterteil wird: die pp-Kette (bei der Protonen zunächst zu Deuterium, dann weiter zu Helium-3 und schließlich Helium-4 verschmelzen) und der CN=-Zyklus, bei dem in Teilreaktionen mit Kohlenstoff, Stickstoff und Sauerstoff Wasserstoff in Helium umgesetzt wird. Hier möchten wir nun eine mögliche Varante des CNO-ZYklus diskutieren.

    \begin{enumerate}
        \item Berechnen Sie hierzu zunächst die \underline{gesamte} freiwerdende Energie über den CNO-Zyklus unter Vernachlässigung des Energieverlusts durch Neutrinos. Beachten Sie dabei die obigen Reaktionsgleichungen sind für die nackten Atomkerne gegeben. Erstellen Sie zunächst die \underline{Gesamtreaktionsbilanz} dieses Zyklus. Erhänzen Sie dann die Gleichung links und rechts mit den notwendigen Elektronenmassen. Berechnen Sie daraus die freiwerdende Energie unter Zuhilfenahme der ATommassen in der NUBASE 2020 der IAEA:
        
        \url{https://www-nds.iaea.org/relnsd/nubase/nubase_min.html}

        \begin{align*}
            ^{12}C + p &\rightarrow ^{13}N + \gamma + Q_1\\
            ^{13}N &\rightarrow ^{13}C + e^+ + \mathbb{v}_e + Q_2\\
        \end{align*}
    \end{enumerate}

\end{document}